\documentclass{article}
\usepackage{amsmath}
\usepackage{amsthm}
\usepackage{amsfonts}
\usepackage{amssymb}
\usepackage{bbold}
\usepackage{dsfont}
\usepackage{graphicx}
\usepackage{cleveref}
\usepackage{bm}
\usepackage{booktabs}
\usepackage{siunitx}
\usepackage{xcolor}
\usepackage{pdfpages}
\usepackage[utf8]{inputenc}
\usepackage[english]{babel}
\newtheorem{definition}{Definition}
\newtheorem{example}{Example}
\newcommand{\exP}[1]{\mathrm{e}^{#1}}
\newcommand{\dif}{\mathrm{d}}
\renewcommand{\vec}[1]{\mathbf{#1}}
\newcommand{\pdagger}{{\phantom{\dagger}}}
\newcommand{\state}[1]{\vert #1\rangle}
\newcommand{\cstate}[1]{\langle #1\vert}
\newcommand{\tr}[1]{\mathrm{tr}\left\{#1\right\}}
\newcommand{\trs}[2]{\, \mathrm{tr}_{#1}\{#2\}}
\newcommand{\expc}[3]{\langle #1\vert #2 \vert #3\rangle}
\usepackage[doi=false,url=false,isbn=false,backend=bibtex,style=nature,natbib=true,sorting=none]{biblatex}
\bibliography{bib.bib}
\usepackage{geometry}
\geometry{
	a4paper,
	total={170mm,257mm},
	left=15mm,
	right=15mm,
	top=15mm,
}
\usepackage{url}
\begin{document}
\title{Literature}
\maketitle
\begin{abstract}
	This is a collection of interesting references, statements, and ideas which I would like to have in my thesis.
\end{abstract}
\section{Ehrenfest}
\subsection{What is missing in Ehrenfest?}
\begin{itemize}
	\item \url{https://pubs.acs.org/doi/full/10.1021/acs.jpclett.3c02690}
	\item Branching is important. 
	\item Decoherence needs to be taken into account properly
	\item Idea: In my approach, I'm not only imposing Ehrenfest, but also single-particle states. Justification (dynamics evolves in a sector with fixed electron number, conservation, but open system??)? 
\end{itemize}
\section{MPS}
\subsection{MPS and bosonic systems}
\begin{itemize}
	\item \url{https://journals.aps.org/pra/abstract/10.1103/PhysRevA.108.042422}
	\item Entanglement area law of ground state holds true even for systems with infinite-dimensional local Hilbert space
	\item Justifies using MPS to describe bosonic systems rigorously
	\item Idea: This is only true for the ground state. But in our transport scenarios, we are not operating in the ground state at all. So the question is still, if MPS are good for far from equilibrium/ground states.
\end{itemize}
\section{Chaos}
\subsection{Chaotic classes in disordered open quantum systems}
\begin{itemize}
	\item \url{https://pubs.aip.org/aip/cha/article/33/11/113113/2920251}
	\item Nomenclature:
	\begin{itemize}
		\item TRS=time-reversal symmetry
		\item SRS=Spin rotation symmetry
		\item CHS=Chiral symmetry
	\end{itemize}
	\item RMT: Tenfold symmetry classification scheme:
	\begin{itemize}
		\item Wigner-Dyson classes (ordinary disordered conductors and chaotic billiards):
		\begin{itemize}
			\item Give rise to an ensemble of Hamiltonian matrices that are randomly distributed according to a probability density distribution $\mathcal{P}(H)$ that are invariant under unitary transformations $H\rightarrow UHU^\dagger$.
			\item Orthogonal: $U^{-1}=U^T$
			\item Unitary: $UU^\dagger = \mathds{1}$
			\item Symplectic, when $U$ are symplectic
		\end{itemize}
		\item Chiral classes (systems that preserve chiral symmetry): 
		\begin{itemize}
			\item Orthogonal (BDI): Preserve TRS, SRS, CHS
			\item Unitary (AIII): TRS is broken (perpendicular magnetic field)
			\item Symplectic (CII): SRS broken, but TRS and CHS remain (spin-orbit coupling)
		\end{itemize}
	\end{itemize}
	\item 
\end{itemize}
\end{document}


